\chapter{Killing 对称、体积与 Hodge 理论}

Levi--Civita 联络把度量与协变微分锁定在同一结构中以后，可以问两个彼此相关的问题。其一，哪些流保持度量，从而代表几何的连续对称？其二，度量如何把微分形式空间变成带内积的 Hilbert 型对象，使外微分 $\dd$ 拥有伴随算子并产生椭圆谱理论？前者引出 Killing 场，后者引出 Hodge 星、余微分与 Hodge 拉普拉斯算子。二者在同一处汇合：Killing 流保持由度量构造的一切自然算子，因此对称性可以与 Hodge 谱同时分解。

\section{Killing 场与测地线守恒量}

设 $(M,g)$ 为 Riemann 或 伪 Riemann 流形，$X$ 为光滑向量场。若其局部流 $\Phi_t$ 满足
\[
\Phi_t^*g=g,
\]
则称 $X$ 为 Killing 向量场。对 $t=0$ 求导，等价地写成
\begin{equation*}\boxed{\mathcal L_Xg=0.}
\end{equation*}
这个定义只表达“流保持度量”，尚未选坐标。使用 Levi--Civita 联络可以把它改写为一阶偏微分方程。

\begin{derivation}[从 Lie 导数得到 Killing 方程，即\eqnrefs{eq:math-dg-synthesis-killing-eq}]
将 $\mathcal L_Xg$ 的定义用 Levi--Civita 联络改写，利用度量相容性和无挠性消去高阶项，得到一阶方程。

对任意向量场 $Y,Z$，Lie 导数定义给出
\begin{equation*}
 (\mathcal L_Xg)(Y,Z)
 =X[g(Y,Z)]-g([X,Y],Z)-g(Y,[X,Z]).
\end{equation*}
这是 Lie 导数的 Cartan 公式在 $(0,2)$ 张量上的特例：第一项是 $X$ 作用于张量值，后两项扣除 $Y,Z$ 自身被 $X$ 的流推动产生的变化。

Levi--Civita 联络满足 $\nabla g=0$，故
\begin{equation*}
 X[g(Y,Z)]
 =(\nabla_Xg)(Y,Z)+g(\nabla_XY,Z)+g(Y,\nabla_XZ)
 =g(\nabla_XY,Z)+g(Y,\nabla_XZ).
\end{equation*}
度量相容性 $\nabla g=0$ 是 Levi--Civita 联络的两条公理之一，几何意义是"协变微分保持内积"。

$\nabla$ 无挠给出 $[X,Y]=\nabla_XY-\nabla_YX$、$[X,Z]=\nabla_XZ-\nabla_ZX$。代入上面的 Lie 导数定义，
\begin{align*}
 (\mathcal L_Xg)(Y,Z)
 &=g(\nabla_XY,Z)+g(Y,\nabla_XZ)\\
 &\quad -g(\nabla_XY-\nabla_YX,Z)-g(Y,\nabla_XZ-\nabla_ZX).
\end{align*}
无挠性是 Levi--Civita 联络的另一条公理，把 Lie 括号这一纯微分拓扑对象用协变导数表达。

含 $\nabla_XY$ 与 $\nabla_XZ$ 的项相消，留下
\begin{equation*}
 (\mathcal L_Xg)(Y,Z)
 =g(\nabla_YX,Z)+g(Y,\nabla_ZX).
\end{equation*}
此式对任意 $Y,Z$ 成立。注意 $X$ 现在被 $\nabla$ 作用，而不是作用于别人——这是 Killing 方程为一阶 PDE 的关键。

取 $Y=\partial_i$、$Z=\partial_j$，$g(\partial_i,\partial_j)=g_{ij}$，$\nabla_{\partial_i}X$ 的 $j$ 分量为 $\nabla_iX^j$，故
\begin{equation*}
 (\mathcal L_Xg)_{ij}
 =g_{jk}\nabla_iX^k+g_{ik}\nabla_jX^k
 =\nabla_iX_j+\nabla_jX_i.
\end{equation*}
$\mathcal L_Xg=0$ 即给出
\begin{equation}\label{eq:math-dg-synthesis-killing-eq}
\boxed{\nabla_iX_j+\nabla_jX_i=0,}
\end{equation}
即\eqnrefs{eq:math-dg-synthesis-killing-eq}。这是一阶线性 PDE， Killing 场构成有限维 Lie 代数，维数至多 $n(n+1)/2$（最大对称空间）。

典型例子。
\begin{itemize}
\item $\RR^n$：平移 $\partial_i$ 与旋转 $x^i\partial_j-x^j\partial_i$ 共 $n+n(n-1)/2=n(n+1)/2$ 个 Killing 场，达到最大维数。
\item $S^n$：$SO(n+1)$ 的 Lie 代数给出 $n(n+1)/2$ 个 Killing 场，对应球面的等距群。
\item $\HH^n$：$SO(n,1)$ 给出 $n(n+1)/2$ 个 Killing 场，双曲空间也是最大对称空间。
\item 一般 Riemann 流形：Killing 场维数 $<n(n+1)/2$，反映对称性的缺失。
\end{itemize}

Killing 场对应度量的连续对称，沿测地线产生守恒量 $J_X=g(X,\dot\gamma)$。在广义相对论中，Killing 场对应时空对称（如稳态时空的时间平移 $\partial_t$ 给出能量守恒、轴对称 $\partial_\phi$ 给出角动量守恒）。Noether 定理在几何语言下即"Killing 场 $\Rightarrow$ 测地线守恒量"。在经典力学中，$n(n+1)/2$ 维最大对称空间对应 $n$ 维自由粒子全部运动积分的存在性。
\end{derivation}

沿仿射参数测地线 $\gamma(s)$，令 $u=\dot\gamma$。Killing 场产生标量
\[
J_X:=g(X,u).
\]
利用 $\nabla_u u=0$ 与\eqnrefs{eq:math-dg-synthesis-killing-eq}，
\[
\begin{aligned}
\frac{\dd J_X}{\dd s}
&=g(\nabla_uX,u)+g(X,\nabla_uu)\\
&=\frac12(\mathcal L_Xg)(u,u)=0.
\end{aligned}
\]
因此
\begin{equation*}\boxed{\frac{\dd}{\dd s}g(X,\dot\gamma)=0.}
\end{equation*}
连续等距对称由此直接转化为测地线的一阶守恒量。

\begin{exbox}[Euclidean 平面的 Killing 代数]
在 $g=\dd x^2+\dd y^2$ 上写
\[
X=P(x,y)\partial_x+Q(x,y)\partial_y.
\]
Killing 方程化为
\[
\partial_xP=0,
\qquad
\partial_yQ=0,
\qquad
\partial_yP+\partial_xQ=0.
\]
故
\[
P=ay+b,
\qquad
Q=-ax+c,
\]
于是全部 Killing 场由
\[
\partial_x,
\qquad
\partial_y,
\qquad
y\partial_x-x\partial_y
\]
张成，分别对应两个平移与一个旋转。
\end{exbox}

\section{体积、Hodge 星与余微分}

若 $(M,g)$ 已定向，度量诱导体积形式。Riemann 情形局部写成
\begin{equation*}\operatorname{vol}_g
=\sqrt{\det(g_{ij})}\,
\dd x^1\wedge\cdots\wedge\dd x^n,
\end{equation*}
对 伪 Riemann 度量则将 $\det g$ 替换为 $|\det g|$。给定任意向量场 $X$，由于 $\mathcal L_X\operatorname{vol}_g$ 仍是顶次形式，存在唯一函数定义散度：
\begin{equation}\label{eq:math-dg-divergence-volume}
\mathcal L_X\operatorname{vol}_g=(\operatorname{div} X)\operatorname{vol}_g.
\end{equation}

由 Cartan 恒等式和顶次形式自动闭合 $\dd\operatorname{vol}_g=0$，
\[
\mathcal L_X\operatorname{vol}_g=\dd(\iota_{X}\operatorname{vol}_g).
\]
在坐标中展开可得
\begin{equation*}\boxed{
\operatorname{div} X
=\nabla_iX^i
=\frac{1}{\sqrt{|\det g|}}
\partial_i\!\left(\sqrt{|\det g|}\,X^i\right).}
\end{equation*}

Killing 方程取迹得到
\[
2\nabla_iX^i=0,
\]
因此 Killing 场满足
\begin{equation*}\operatorname{div} X=0,
\qquad
\mathcal L_X\operatorname{vol}_g=0.
\end{equation*}
等距流不仅保持长度和角度，也保持由同一度量诱导的体积。

从体积形式继续到微分形式的度量对偶。Hodge 星把同一度量对体积与内积的作用统一起来，并自然导出余微分与 Hodge 拉普拉斯算子。

从这里起假设 $(M,g)$ 为定向 Riemann 流形。度量在 $k$-形式上诱导点态内积 $\langle\cdot,\cdot\rangle_g$。Hodge 星算子
\begin{equation*}*:\Omega^k(M)\to\Omega^{n-k}(M)
\end{equation*}
由
\begin{equation}\label{eq:math-dg-hodge-star-characterization}
\alpha\wedge *\beta
=\langle\alpha,\beta\rangle_g\,\operatorname{vol}_g
\end{equation}
唯一确定。对 Riemann 度量，
\[
**\alpha=(-1)^{k(n-k)}\alpha.
\]
因此 $*$ 不是额外选择，而是“度量 + 定向”在外代数上的自然延伸。

在紧支撑形式上定义 $L^2$ 内积
\[
(\alpha,\beta)_{L^2}
:=\int_M\alpha\wedge *\beta.
\]
外微分的形式伴随称为余微分。对 $k$-形式采用约定
\begin{equation*}\delta
:=(-1)^{n(k+1)+1}*\dd*:
\Omega^k(M)\to\Omega^{k-1}(M).
\end{equation*}
若流形无边界，或形式满足足以消去边界项的条件，则 Stokes 定理给出
\[
(\dd\alpha,\beta)_{L^2}
=(\alpha,\delta\beta)_{L^2}.
\]
特别地，对向量场 $X$ 的对偶一形式 $X^\flat$，
\begin{equation*}\boxed{\operatorname{div} X=-\delta X^\flat=*\dd*X^\flat.}
\end{equation*}
最后一个等号使用的是 $k=1$ 时 $\delta=-*\dd*$。这一定号关系必须与\eqnrefs{eq:math-dg-divergence-volume}一致；若把 $-*\dd*X^\flat$ 直接写成散度，会与当前余微分约定差一个负号。

Hodge--de Rham 拉普拉斯算子 定义为
\begin{equation*}\Delta:=\dd\delta+\delta\dd.
\end{equation*}
它保持形式次数。对任意紧支撑 $k$-形式，
\begin{equation}\label{eq:math-dg-hodge-energy}
(\alpha,\Delta\alpha)_{L^2}
=\|\dd\alpha\|_{L^2}^2+\|\delta\alpha\|_{L^2}^2\ge0.
\end{equation}
因此在无边界紧致 Riemann 流形上，$\Delta$ 是非负自伴椭圆算子；这一点把本章与前面的谱理论直接接起来。

\begin{derivation}[Hodge 拉普拉斯算子 的主符号与椭圆性]
计算 $\Delta=\dd\delta+\delta\dd$ 的二阶主符号。$\dd$ 与 $\delta$ 的一阶主符号分别由外积与插入算子给出，二者复合后用 Cartan 恒等式 $\iota_{\xi^\sharp}(\xi\wedge\cdot)+\xi\wedge\iota_{\xi^\sharp}\cdot=|\xi|^2\,\cdot$ 化为 $|\xi|^2\id()$。

在点 $p$ 处取非零余向量 $\xi\in T_p^*M$。采用 $\partial_j\mapsto \ii\xi_j$ 的主符号约定，外微分与余微分的一阶主符号分别为
\begin{equation*}
\sigma_1(\dd)(\xi)=\ii\,\xi\wedge,
\qquad
\sigma_1(\delta)(\xi)=-\ii\,\iota_{\xi^\sharp}.
\end{equation*}
$\sigma_1(\dd)=\ii\,\xi\wedge$ 来自外微分对一形式 $\alpha$ 给 $\dd\alpha=\partial_i\alpha_j\,\dd x^i\wedge\dd x^j$，主项含一阶导数。$\sigma_1(\delta)=-\ii\,\iota_{\xi^\sharp}$ 来自余微分的 Hodge 星定义与 $\delta=\pm *\dd*$ 中 $*$ 不改变阶数。

因此二阶算子 $\Delta=\dd\delta+\delta\dd$ 的主符号为
\begin{equation*}
\begin{aligned}
\sigma_2(\Delta)(\xi)
&=\sigma_1(\dd)\sigma_1(\delta)+\sigma_1(\delta)\sigma_1(\dd)\\
&=(\ii\,\xi\wedge)(-\ii\,\iota_{\xi^\sharp})
+(-\ii\,\iota_{\xi^\sharp})(\ii\,\xi\wedge)\\
&=(\xi\wedge)\iota_{\xi^\sharp}
+\iota_{\xi^\sharp}(\xi\wedge).
\end{aligned}
\end{equation*}
两个 $\ii$ 与 $-\ii$ 相乘都给 $+1$，符号约定使最终结果为正。

对任意形式 $\alpha$，插入算子与外积满足
\begin{equation*}
\iota_{\xi^\sharp}(\xi\wedge\alpha)
=|\xi|_g^2\alpha
-\xi\wedge\iota_{\xi^\sharp}\alpha.
\end{equation*}
这是 Cartan 公式 $\mathcal L_{\xi^\sharp}=\dd\iota_{\xi^\sharp}+\iota_{\xi^\sharp}\dd$ 在常向量场上的特例，也可直接对单项式 $\alpha=\theta^{i_1}\wedge\cdots\wedge\theta^{i_k}$ 验证：若 $\xi$ 出现在 $\alpha$ 中，插入与外积相继给出 $|\xi|^2$；若不出现，外积先添加 $\xi$ 再插入移除，给零。

把 Cartan 恒等式代入主符号表达式，
\begin{equation*}
\iota_{\xi^\sharp}(\xi\wedge\alpha)
+\xi\wedge\iota_{\xi^\sharp}\alpha
=|\xi|_g^2\alpha
-\xi\wedge\iota_{\xi^\sharp}\alpha
+\xi\wedge\iota_{\xi^\sharp}\alpha
=|\xi|_g^2\alpha,
\end{equation*}
两项相消，恰好得到
\begin{equation}\label{eq:math-dg-hodge-principal-symbol}
\boxed{
\sigma_2(\Delta)(\xi)
=|\xi|_g^2\id(\Lambda^kT_p^*M).}
\end{equation}
当 $\xi\neq0$ 时该线性映射可逆，故 $\Delta$ 是椭圆算子。结合\eqnrefs{eq:math-dg-hodge-energy}的非负性与适当自伴定义域，就能调用前面建立的自伴谱理论，而不必把 Hodge 分解当作孤立结论。

典型例子与对比。
\begin{itemize}
\item $\RR^n$ 上 $0$-形式（函数）：$\Delta f=-\sum_i\partial_i^2 f$（本书非负约定），主符号 $|\xi|^2$，正是普通 拉普拉斯算子。
\item $\RR^n$ 上 $1$-形式：$\Delta\alpha_j=-\sum_i\partial_i^2\alpha_j$，每个分量独立满足标量 拉普拉斯算子，主符号仍是 $|\xi|^2\id()$。
\item 一般流形：主符号与 $\RR^n$ 相同，但低阶项含曲率（见 Weitzenböck 公式 $\Delta_H=\nabla^*\nabla+\operatorname{Ric}$），曲率通过低阶项影响谱而非椭圆性。
\end{itemize}

Hodge 拉普拉斯算子 的椭圆性是 Hodge 定理成立的分析基础：椭圆 + 自伴 + 非负 + 紧致流形 $\Rightarrow$ 谱离散且有限维核。在物理中，这对应规范理论中"规范场方程的解空间有限维"（如 Yang--Mills 瞬子模空间），以及弦理论中 拉普拉斯算子 谱决定物理态能级。Weitzenböck 公式把椭圆性 + 曲率正性结合，给出 Bochner 消没定理——曲率通过低阶项直接控制核空间维数。
\end{derivation}

Hodge 拉普拉斯算子 的主符号恰为 $|\xi|_g^2$（\eqnrefs{eq:math-dg-hodge-principal-symbol}），这是其椭圆性与谱离散性的来源。

\section{Hodge 分解与上同调}

记
\[
\mathcal H^k(M):=\ker\Delta\big|_{\Omega^k(M)}
\]
为 调和 $k$-形式空间。由\eqnrefs{eq:math-dg-hodge-energy}，
\begin{equation}\label{eq:math-dg-harmonic-closed-coclosed}
\Delta\alpha=0
\quad\Longleftrightarrow\quad
\dd\alpha=0,
\ \delta\alpha=0.
\end{equation}

\begin{theorem}[紧致无边界流形上的 Hodge 分解]\label{thm:math-dg-hodge-decomposition}
设 $(M,g)$ 为紧致、定向、无边界 Riemann 流形。则任意 $\alpha\in\Omega^k(M)$ 都有正交分解
\begin{equation}\label{eq:math-dg-hodge-decomposition}
\boxed{
\alpha=\dd\beta+\delta\gamma+\alpha_{\mathcal H},
\qquad
\alpha_{\mathcal H}\in\mathcal H^k(M).}
\end{equation}
其中 精确、coexact 与 调和 三个分量作为形式是唯一的；势 $\beta$ 与 $\gamma$ 本身一般不唯一。并且自然映射
\begin{equation*}\mathcal H^k(M)\longrightarrow H_{\mathrm{dR}}^k(M),
\qquad
\eta\longmapsto[\eta]
\end{equation*}
是线性同构。
\end{theorem}

\begin{derivation}[Green 算子给出的 Hodge 分解]
在紧致无边界流形上，限制在 $\Omega^k(M)$ 的椭圆自伴算子记为 $\Delta_k$。它具有离散非负谱，零特征空间就是有限维的 $\mathcal H^k(M)$。记到该零空间的正交投影为 $\Pi_{\mathcal H,k}$，并在其正交补上定义 Green 算子
\[
\mathcal G_k:=\left(\Delta_k\big|_{\mathcal H^k(M)^\perp}\right)^{-1},
\qquad
\mathcal G_k\eta=0\quad(\eta\in\mathcal H^k(M)).
\]
于是谱分解给出精确算子恒等式
\[
I=\Pi_{\mathcal H,k}+\Delta_k\mathcal G_k.
\]
对任意 $\alpha\in\Omega^k(M)$，直接展开 $\Delta_k=\dd\delta+\delta\dd$ 得
\[
\begin{aligned}
\alpha
&=\Pi_{\mathcal H,k}\alpha+\Delta_k\mathcal G_k\alpha\\
&=\Pi_{\mathcal H,k}\alpha
+\dd\delta\mathcal G_k\alpha
+\delta\dd\mathcal G_k\alpha.
\end{aligned}
\]
令
\[
\alpha_{\mathcal H}=\Pi_{\mathcal H,k}\alpha,
\qquad
\beta=\delta\mathcal G_k\alpha,
\qquad
\gamma=\dd\mathcal G_k\alpha,
\]
便得到\eqnrefs{eq:math-dg-hodge-decomposition}。

三个分量两两正交。例如
\[
(\dd\beta,\delta\gamma)
=(\dd^2\beta,\gamma)=0,
\]
而对 调和 分量 $\alpha_{\mathcal H}$，由 $\dd\alpha_{\mathcal H}=\delta\alpha_{\mathcal H}=0$，
\[
(\dd\beta,\alpha_{\mathcal H})=(\beta,\delta\alpha_{\mathcal H})=0,
\qquad
(\delta\gamma,\alpha_{\mathcal H})=(\gamma,\dd\alpha_{\mathcal H})=0.
\]
因此三项本身唯一。若只改变 $\beta$ 为 $\beta+\zeta$ 且 $\dd\zeta=0$，精确 分量不变，所以势并不唯一；$\gamma$ 同理。

最后证明调和代表与 de Rham 上同调一一对应。调和形式由\eqnrefs{eq:math-dg-harmonic-closed-coclosed}必为闭形式，故定义上同调类。若 $\eta_1-\eta_2=\dd\zeta$ 且 $\eta_1,\eta_2$ 都 调和，则
\[
\|\eta_1-\eta_2\|^2
=(\eta_1-\eta_2,\dd\zeta)
=(\delta(\eta_1-\eta_2),\zeta)=0,
\]
所以同一上同调类至多一个 调和 代表。另一方面若 $\alpha$ 闭，Hodge 分解给出
\[
\alpha=\dd\beta+\delta\gamma+\alpha_{\mathcal H}.
\]
对两边作用 $\dd$，得到 $\dd\delta\gamma=0$。于是
\[
\|\delta\gamma\|^2
=(\gamma,\dd\delta\gamma)=0,
\]
故 $\delta\gamma=0$，从而 $[\alpha]=[\alpha_{\mathcal H}]$。存在性成立。
\end{derivation}

意义：Hodge 分解 $\Omega^k=\dd\Omega^{k-1}\oplus\delta\Omega^{k+1}\oplus\mathcal H^k$ 把任意微分形式正交分解为恰当、余恰当和调和三部分。每个 de Rham 上同调类有且仅有一个调和代表元——这使得上同调（拓扑不变量）可用调和形式（分析对象）实现。实例：紧曲面 $\Sigma_g$ 上，$H^1_{\mathrm{dR}}$ 由 $2g$ 个线性无关的调和 1-形式张成，对应 $g$ 个环柄的经线和纬线方向。在物理中，Maxwell 方程 $\dd F=0$、$\delta F=J$ 的解存在性恰好由 Hodge 分解保证：$F=F_{\mathcal H}+\dd A+\delta B$，规范势 $A$ 存在当且仅当 $[F]\in H^2$ 被源 $J$ 的约束允许。



Hodge 定理还把一个纯拓扑结论变成了几乎一行的 $L^2$ 论证。设 $M$ 为紧致、定向、无边界 $n$ 维流形。Hodge 星给出
\[
*:\Omega^k(M)\longrightarrow\Omega^{n-k}(M),
\]
并与 Hodge 拉普拉斯算子 对易：
\[
\Delta_H(*\alpha)=*(\Delta_H\alpha).
\]
因此 $*$ 限制为 调和 空间之间的同构
\begin{equation}
*:\mathcal H^k(M)\overset{\sim}{\longrightarrow}\mathcal H^{n-k}(M).
\label{eq:math-dg-hodge-star-harmonic-isomorphism}
\end{equation}
结合 Hodge--de Rham 同构，就得到 Poincaré 对偶 的解析实现。

\begin{theorem}[Poincaré 对偶的 Hodge 实现]\label{thm:math-dg-poincare-duality-hodge}
对紧致、定向、无边界 $n$ 维流形 $M$，配对
\begin{equation}
H^k_{\mathrm{dR}}(M)\times H^{n-k}_{\mathrm{dR}}(M)
\longrightarrow\RR,
\qquad
([\alpha],[\beta])\longmapsto\int_M\alpha\wedge\beta
\label{eq:math-dg-poincare-pairing}
\end{equation}
是非退化的。特别地，
\[
b_k(M)=b_{n-k}(M).
\]
\end{theorem}

\begin{derivation}[从调和代表元到 Poincaré 对偶]
先验证\eqnrefs{eq:math-dg-poincare-pairing}只依赖上同调类。若 $\alpha\mapsto\alpha+\dd\eta$ 且 $\beta$ 闭，则由 $\dd\beta=0$ 及外微分 Leibniz 律
\[
\dd(\eta\wedge\beta)
=\dd\eta\wedge\beta+(-1)^{\deg\eta}\eta\wedge\dd\beta
=\dd\eta\wedge\beta,
\]
故由 Stokes 定理及 $\partial M=\varnothing$，
\[
\int_M\dd\eta\wedge\beta
=\int_M\dd(\eta\wedge\beta)
=\int_{\partial M}\eta\wedge\beta
=0.
\]
对 $\beta$ 的 精确 改变 $\beta\mapsto\beta+\dd\zeta$ 同理，因 $\alpha$ 闭且
\[
\dd(\alpha\wedge\zeta)
=\dd\alpha\wedge\zeta+(-1)^{\deg\alpha}\alpha\wedge\dd\zeta
=\alpha\wedge\dd\zeta,
\]
故
\[
\int_M\alpha\wedge\dd\zeta
=\int_M\dd(\alpha\wedge\zeta)=0.
\]
因此配对在上同调上良定。

现在取非零上同调类 $[\alpha]\neq0$，并令 $h$ 为其唯一调和代表元，即 $\Delta h=0$。由\eqnrefs{eq:math-dg-hodge-star-harmonic-isomorphism}，Hodge 星保持 调和 性，故 $*h$ 也是 调和，因此定义一个 $(n-k)$ 次上同调类
\[
[*h]\in H^{n-k}_{\mathrm{dR}}(M).
\]
于是由 Hodge 星的定义\eqnrefs{eq:math-dg-hodge-star-characterization}，
\begin{align*}
\int_M h\wedge *h
&=\int_M \langle h,h\rangle_g\,\operatorname{vol}_g\\
&=\int_M |h|_g^2\,\operatorname{vol}_g\\
&=\|h\|_{L^2}^2.
\end{align*}
由 $[\alpha]\neq0$ 故 $h\neq0$，从而
\[
\|h\|_{L^2}^2>0.
\]
所以 $[\alpha]$ 与 $[*h]$ 配对非零，不可能与所有 $H^{n-k}_{\mathrm{dR}}(M)$ 元素配对为零。反向非退化性由 $**h=(-1)^{k(n-k)}h$ 完全相同地得到：对非零 $[\beta]\in H^{n-k}_{\mathrm{dR}}(M)$，取其 调和 代表元 $k$，则
\[
\int_M k\wedge *k=\|k\|_{L^2}^2>0,
\]
故 $[\beta]$ 与 $[*k]\in H^k_{\mathrm{dR}}(M)$ 配对非零。因此配对非退化，而有限维向量空间之间的非退化配对立刻给出
\[
\dim H^k_{\mathrm{dR}}(M)
=\dim H^{n-k}_{\mathrm{dR}}(M).
\]

紧曲面 $\Sigma_g$（亏格 $g$）满足 $b_0=b_2=1$、$b_1=2g$，因而 $b_k=b_{2-k}$ 逐次成立；中间次数 $k=1$ 是自对偶的。紧 3-流形如 $T^3$ 则有 $b_0=b_3=1$、$b_1=b_2=3$。Poincaré 对偶不仅是维数等式，还给出上积配对 $H^k\times H^{n-k}\to H^n\cong\mathbb R$ 的非退化性——这是相交理论的基础：$k$-形式与 $(n-k)$-形式的上积对应对偶子流形的相交数。
\end{derivation}

这个证明说明定向、紧致和无边界三个条件分别在何处进入：定向保证可全局积分顶次形式；紧致保证积分与 调和 空间的有限维谱理论良好；无边界使 Stokes 边界项消失。若有边界，正确对象会变成 绝对/相对 上同调之间的 Lefschetz 对偶，正好与后面的两类 Hodge 边界条件对应。
Hodge 拉普拉斯算子 的定义使用外微分与余微分，而 Riemann 几何中还存在由 Levi--Civita 联络构造的 rough 拉普拉斯算子
\[
\nabla^*\nabla:
\Omega^k(M)\to\Omega^k(M).
\]
二者具有相同的二阶主符号，因此差只能是零阶曲率作用。一般 Weitzenböck 公式写成
\begin{equation*}\boxed{
\Delta_H
=\nabla^*\nabla+\mathcal R_k
},\end{equation*}
其中 $\mathcal R_k$ 是 Riemann 曲率在 $k$-形式表示上的代数作用。对 $0$-形式，$\mathcal R_0=0$；对 $1$-形式最简单：
\begin{equation}
\boxed{
(\Delta_H\alpha)_i
=(\nabla^*\nabla\alpha)_i
+\operatorname{Ric}_i{}^j\alpha_j
}.
\label{eq:math-dg-weitzenbock-oneform}
\end{equation}
这条公式把拓扑信息、椭圆 PDE 与 Ricci 曲率直接连接起来。

若 $M$ 紧致无边界且 $\alpha$ 为 调和 $1$-形式，则
\[
0=(\alpha,\Delta_H\alpha)
=\|\nabla\alpha\|_{L^2}^2
+\int_M\operatorname{Ric}(\alpha^\sharp,\alpha^\sharp)\,\operatorname{vol}_g.
\]
若 $\operatorname{Ric}>0$ 处处正定，两项都非负且第二项只有 $\alpha=0$ 才能消失，因此
\[
\mathcal H^1(M)=0,
\qquad
H^1_{\mathrm{dR}}(M)=0.
\]
这就是最基本的 Bochner vanishing 定理：曲率的正性通过 Weitzenböck 零阶项排除 调和 代表元，从而杀掉上同调。

\begin{derivation}[一形式 Weitzenböck 公式的指标推导]
对 $1$-形式 $\alpha$，分别计算 $\delta\dd\alpha$ 与 $\dd\delta\alpha$ 的分量，相加后把"协变导数交换子"识别为 Ricci 曲率，得到 $\Delta_H=\nabla^*\nabla+\operatorname{Ric}$。

对 $1$-形式 $\alpha=\alpha_j\dd x^j$，采用
\begin{equation*}
(\dd\alpha)_{ij}
=\nabla_i\alpha_j-\nabla_j\alpha_i,
\qquad
\delta\alpha=-\nabla^i\alpha_i.
\end{equation*}
第一个是外微分的协变分量形式（无挠使 Christoffel 项相消），第二个是余微分的分量约定（散度的负号来自本书非负 拉普拉斯算子 约定）。

把 $\dd\alpha$ 视为 $2$-形式再作用 $\delta$：
\begin{equation*}
(\delta\dd\alpha)_j
=-\nabla^i(\dd\alpha)_{ij}
=-\nabla^i(\nabla_i\alpha_j-\nabla_j\alpha_i)
=-\nabla^i\nabla_i\alpha_j+\nabla^i\nabla_j\alpha_i.
\end{equation*}
第一项是粗 拉普拉斯算子（联络 拉普拉斯算子，亦称 $\nabla^*\nabla$ 的负），第二项保留待与 $\dd\delta\alpha$ 中的对应项合并。


\begin{equation*}
(\dd\delta\alpha)_j
=\nabla_j(\delta\alpha)
=-\nabla_j\nabla^i\alpha_i.
\end{equation*}
此处用了 $1$-形式 $\delta\alpha$ 的外微分分量公式 $(\dd f)_j=\nabla_j f$。

相加得到
\begin{equation*}
(\Delta_H\alpha)_j
=(\delta\dd\alpha+\dd\delta\alpha)_j
=-\nabla^i\nabla_i\alpha_j
+\nabla^i\nabla_j\alpha_i
-\nabla_j\nabla^i\alpha_i.
\end{equation*}
最后两项正是协变导数交换子。对协变 $1$-张量，
\begin{equation*}
(\nabla^i\nabla_j-\nabla_j\nabla^i)\alpha_i
=\operatorname{Ric}_j{}^k\alpha_k
\end{equation*}
（符号与本书前面的 Riemann 曲率 约定 相容）。这一识别来自 Riemann 曲率的分量定义 $[\nabla_a,\nabla_b]X^c=R^c{}_{dab}X^d$，对协变指标取 迹 给 Ricci：$R^i{}_{j i k}\alpha^k=\operatorname{Ric}_{jk}\alpha^k$。

另一方面，正规坐标中
\begin{equation*}
(\nabla^*\nabla\alpha)_j
=-\nabla^i\nabla_i\alpha_j
\end{equation*}
在基点成立，张量性使其全局成立，因此得到\eqnrefs{eq:math-dg-weitzenbock-oneform}。

典型例子与应用。
\begin{itemize}
\item $\RR^n$：$\operatorname{Ric}=0$，$\Delta_H=\nabla^*\nabla=-\sum_i\partial_i^2$，每个分量独立满足标量 拉普拉斯算子。
\item $S^n$（$n\ge2$，Ric $=(n-1)g$）：$\Delta_H=\nabla^*\nabla+(n-1)$。调和 $1$-形式满足 $\nabla^*\nabla\alpha=-(n-1)\alpha$，但 $\nabla^*\nabla$ 非负，故 $\alpha=0$，即 $H^1(S^n)=0$（$n\ge2$）。
\item 紧致流形 $\operatorname{Ric}>0$：Bochner 消没定理给 $H^1(M)=0$，结合 Bonnet--Myers 定理（$\operatorname{Ric}\ge(n-1)k>0$ $\Rightarrow$ 直径 $\le\pi/\sqrt{k}$）说明正 Ricci 曲率强烈限制流形拓扑。
\end{itemize}

更一般结论——Lichnerowicz 定理。对 Dirac 算子 $D$，类似推导给 Lichnerowicz 公式 $D^2=\nabla^*\nabla+\frac14\,\mathrm{Scal}$。在紧致流形 $\mathrm{Scal}>0$ 上，$D^2$ 严格正，故 $\ker D=0$，结合 Atiyah--Singer 指标定理 $\operatorname{ind}(D)=\hat A(M)$ 给 $\hat A(M)=0$。这是"曲率正性 $\Rightarrow$ 拓扑消没"的典型链条，热核方法通过 $\operatorname{Str}(e^{-tD^2})$ 把这一局部曲率信息与全局指标精确关联。
\end{derivation}


Bochner 方法还给出一个比“$H^1$ 消失”更定量的结论。对函数 $f$，采用本书非负 拉普拉斯算子 约定
\[
\Delta f=-\nabla^i\nabla_i f,
\]
Bochner 恒等式可写成
\begin{equation}
\frac12\Delta|\nabla f|^2
=\langle\nabla f,\nabla\Delta f\rangle
-|\nabla^2f|^2
-\operatorname{Ric}(\nabla f,\nabla f).
\label{eq:math-dg-bochner-function}
\end{equation}
若 $M$ 紧致且 $\operatorname{Ric}\ge(n-1)k g$，$k>0$，对第一非零本征函数
\[
\Delta f=\lambda_1 f,
\qquad
\int_M f\,\operatorname{vol}_g=0,
\]
积分\eqnrefs{eq:math-dg-bochner-function}并用
\[
|\nabla^2 f|^2\ge\frac1n(\tr\nabla^2f)^2
=\frac1n(\Delta f)^2
\]
可得
\begin{equation}
\boxed{\lambda_1\ge nk.}
\label{eq:math-dg-lichnerowicz-gap}
\end{equation}
这就是 Lichnerowicz 谱 隙。它与前一章 Bonnet--Myers 的假设完全相同，但结论从“直径有限”变成“拉普拉斯算子 的第一个正本征值远离零”。同一个 Ricci 下界同时约束几何尺度与低频谱。

\begin{derivation}[从 Bochner 恒等式到 Lichnerowicz 下界，即\eqnrefs{eq:math-dg-lichnerowicz-gap}]
设 $(M^n,g)$ 紧致无边界，$\operatorname{Ric}\ge(n-1)kg$（$k>0$），$\Delta f=\lambda_1 f$（第一非零本征函数）。

对本征函数积分\eqnrefs{eq:math-dg-bochner-function}。紧致无边界时 Stokes 定理给出 $\int_M\Delta|\nabla f|^2\operatorname{vol}_g=0$（边界项消失），于是
\begin{equation*}
 0
 =\lambda_1\int_M|\nabla f|^2\operatorname{vol}_g
 -\int_M|\nabla^2f|^2\operatorname{vol}_g
 -\int_M\operatorname{Ric}(\nabla f,\nabla f)\operatorname{vol}_g.
\end{equation*}

由 Green 公式，
\begin{equation*}
 \int_M|\nabla f|^2\operatorname{vol}_g
 =\int_M f\,\Delta f\,\operatorname{vol}_g
 =\lambda_1\int_M f^2\operatorname{vol}_g.
\end{equation*}

由 Cauchy--Schwarz 不等式对 Hessian 的迹，
\begin{equation*}
 |\nabla^2 f|^2\ge\frac{1}{n}(\tr\nabla^2f)^2
 =\frac{1}{n}(\Delta f)^2
 =\frac{\lambda_1^2}{n}f^2.
\end{equation*}
由 Ricci 下界假设，
\begin{equation*}
 \operatorname{Ric}(\nabla f,\nabla f)\ge(n-1)k\,|\nabla f|^2
 =(n-1)k\lambda_1\,f^2.
\end{equation*}

将两个下界代入上面的 Bochner 恒等式（移项后右端 $\ge0$）：
\begin{align*}
 0
 &\le \lambda_1^2\int_Mf^2\operatorname{vol}_g
  -\frac{\lambda_1^2}{n}\int_Mf^2\operatorname{vol}_g
  -(n-1)k\lambda_1\int_Mf^2\operatorname{vol}_g\\
 &=\lambda_1\left(\frac{(n-1)\lambda_1}{n}-(n-1)k\right)\int_Mf^2\operatorname{vol}_g\\
 &=(n-1)\lambda_1\left(\frac{\lambda_1}{n}-k\right)\int_Mf^2\operatorname{vol}_g.
\end{align*}

因为 $\lambda_1>0$（非常数本征函数）、$n>1$、$f\not\equiv0$，积分因子为正，故 $\lambda_1/n-k\ge0$，即 $\lambda_1\ge nk$。这就是\eqnrefs{eq:math-dg-lichnerowicz-gap}。

Lichnerowicz 下界与 Bonnet--Myers 直径上界 $\operatorname{diam}(M)\le\pi/\sqrt{k}$ 假设完全相同（$\operatorname{Ric}\ge(n-1)kg$），但结论从"直径有限"变成"拉普拉斯算子 的第一个正本征值远离零"。同一个 Ricci 下界同时约束几何尺度与低频谱。
\end{derivation}
Hodge 分解之所以能从 $L^2$ 正交分解提升到光滑微分形式分解，关键不是纯代数，而是椭圆正则性。对紧致无边界 $M$，任意实数 $s$ 都有估计
\begin{equation}
\boxed{
\|\alpha\|_{H^{s+2}}
\le C_s\bigl(
\|\Delta_H\alpha\|_{H^s}
+\|\alpha\|_{L^2}
\bigr)
}.
\label{eq:math-dg-hodge-elliptic-estimate}
\end{equation}
因此如果 $\Delta_H\alpha$ 比 $\alpha$ 已知得更光滑两阶，$\alpha$ 本身就自动获得两阶正则性。特别地，若 $\Delta_H\alpha=0$ 且 $\alpha\in L^2$，反复 自助法 得
\[
\alpha\in H^s\quad\forall s,
\]
再由 Sobolev 嵌入 得 $\alpha\in C^\infty$。所以 调和形式 的光滑性不是定义里额外要求的，而是椭圆方程自动推出的。

把\eqnrefs{eq:math-dg-hodge-elliptic-estimate}与 Rellich 紧嵌入
\[
H^{s+2}(M)\hookrightarrow H^s(M)
\]
结合，可知 预解式 $(\Delta_H+I)^{-1}:L^2\to L^2$ 是紧算子。于是前面的紧自伴谱理论立即给出
\[
0\le\lambda_0\le\lambda_1\le\cdots,
\qquad
\lambda_j\to\infty,
\]
每个本征值有限重，且本征形式构成 $L^2$ 正交完备系。调和形式 正是 $\lambda=0$ 的有限维本征空间。

更细的高能渐近由 Weyl law 控制。对 $k$-形式 拉普拉斯算子，若 $N_k(\Lambda)$ 表示不超过 $\Lambda$ 的本征值计数（计重数），则主项只由主符号和体积决定：
\begin{equation*}N_k(\Lambda)
\sim
\binom{n}{k}
\frac{\omega_n}{(2\pi)^n}
\operatorname{Vol}_g(M)
\Lambda^{n/2},
\qquad \Lambda\to\infty.\end{equation*}
曲率影响低阶修正，而最高阶密度与 Euclidean 拉普拉斯算子 相同。这再次体现 椭圆 算符 的局部主符号决定高频、全局几何决定低频零模与谱细节。

Green 算子 $\mathcal G_k$ 因此不仅是 Hodge 分解中的形式逆，还具有两阶 光滑化：在 调和 正交 补 上
\[
\mathcal G_k:H^s\to H^{s+2}
\]
连续。若写成谱展开
\[
\mathcal G_k\alpha
=\sum_{\lambda_j>0}
\frac{\langle\alpha,\phi_j\rangle}{\lambda_j}\phi_j,
\]
就能直接看出高频系数被 $1/\lambda_j$ 抑制。这与前面逆问题章节的 谱 滤波器 完全同构；区别只在于 Hodge 拉普拉斯算子 在紧流形上具有离散、非负、自伴谱，并且零模由拓扑决定。


同一个谱分解还自然生成 热 semigroup
\begin{equation*}\mathrm e^{-t\Delta_H}\alpha
=\sum_{j=0}^{\infty}
\mathrm e^{-t\lambda_j}
\langle\alpha,\phi_j\rangle\phi_j,
\qquad t\ge0.\end{equation*}
它解决初值问题
\[
\partial_tu+\Delta_Hu=0,
\qquad
u(0)=\alpha.
\]
对任意 $t>0$，指数因子比任何多项式更快压低高频，因此 $\mathrm e^{-t\Delta_H}$ 是无穷阶 光滑化 算子；而 $t\to\infty$ 时只有零本征空间留下：
\begin{equation*}\mathrm e^{-t\Delta_H}\alpha
\longrightarrow
\Pi_{\mathcal H}\alpha
\quad\text{in }L^2.\end{equation*}
因此 调和 投影 也可理解为“把任意微分形式在 Hodge 热 流 下演化到无限久”的极限。

Green 算子则是同一半群的时间积分：在 $\mathcal H^k(M)^\perp$ 上，
\begin{equation*}\boxed{
\mathcal G_k
=\int_0^{\infty}\mathrm e^{-t\Delta_k}\,\dd t,
}\end{equation*}
而在全空间上更准确地写成
\[
\mathcal G_k
=\int_0^{\infty}
\left(\mathrm e^{-t\Delta_k}-\Pi_{\mathcal H,k}\right)\dd t.
\]
逐个本征模积分
\[
\int_0^\infty\mathrm e^{-t\lambda_j}\dd t=\frac1{\lambda_j}
\qquad(\lambda_j>0)
\]
便恢复 Green 谱展开。于是 Hodge 分解、热方程、Green 算子和拓扑零模其实是同一个非负自伴算子的四种表示。

有边界时，Hodge 分解还需要把 Green 公式留下的边界型纳入算子定义域；绝对/相对 条件正是这一步的两个自然自伴实现。

无边界情形中，Stokes 定理没有留下边界项，因此 $\dd$ 与 $\delta$ 自动互为形式伴随；一旦 $\partial M\neq\varnothing$，这个结论必须由边界条件补回。设 $M$ 为紧致定向 Riemann 流形，$\nu$ 是外法单位向量，$\nu^\flat$ 是其度量对偶一形式，$j:\partial M\hookrightarrow M$ 为包含映射。对 $\alpha\in\Omega^{k-1}(M)$、$\beta\in\Omega^k(M)$，Green 公式为
\begin{equation}\label{eq:math-dg-hodge-green-boundary}
 (\dd\alpha,\beta)_{L^2}-(\alpha,\delta\beta)_{L^2}
 =\int_{\partial M}
 \left\langle \nu^\flat\wedge\alpha,\beta\right\rangle_g\dd S
 =\int_{\partial M}
 \left\langle \alpha,\iota_{\nu}\beta\right\rangle_g\dd S.
\end{equation}
因此“调和”在有边界情形不能只写 $\Delta\alpha=0$；必须同时指定使 $\Delta$ 成为自伴椭圆算子的边界域。

对 $k$-形式 $\omega$，绝对（又称 Neumann 型）边界条件定义为
\begin{equation*}\boxed{
 \iota_{\nu}\omega\big|_{\partial M}=0,
 \qquad
 \iota_{\nu}\dd\omega\big|_{\partial M}=0.}
\end{equation*}
第一式消去 $\omega$ 的法向分量，第二式是与 拉普拉斯算子 二阶性配对的自然补充条件。相对（又称 Dirichlet 型）边界条件定义为
\begin{equation*}\boxed{
 j^*\omega=0,
 \qquad
 j^*(\delta\omega)=0.}
\end{equation*}
由于在边界上任意形式唯一分解成切向部分与 $\nu^\flat$ 楔法向部分，第一条 相对 条件也等价于
\[
 \nu^\flat\wedge\omega\big|_{\partial M}=0.
\]
Hodge 星交换这两类边界条件：若 $\omega$ 满足 绝对 条件，则 $*\omega$ 满足 相对 条件，反之亦然。

记相应自伴实现为 $\Delta_{A,k}$ 与 $\Delta_{R,k}$，并定义
\begin{equation*}\mathcal H_A^k(M):=\ker\Delta_{A,k},
 \qquad
 \mathcal H_R^k(M):=\ker\Delta_{R,k}.
\end{equation*}
在任一边界域上，Green 公式的边界项都消失，从而能量恒等式仍成立：
\begin{equation}\label{eq:math-dg-boundary-hodge-energy}
 (\omega,\Delta_{A/R,k}\omega)_{L^2}
 =\|\dd\omega\|_{L^2}^2+\|\delta\omega\|_{L^2}^2.
\end{equation}
故 绝对/相对 调和形式 都同时闭且余闭。真正改变的是它们代表的拓扑对象：Hodge--Morrey--Friedrichs 定理给出自然同构
\begin{equation*}\boxed{
 \mathcal H_A^k(M)\simeq H_{\mathrm{dR}}^k(M),
 \qquad
 \mathcal H_R^k(M)\simeq H_{\mathrm{dR}}^k(M,\partial M).}
\end{equation*}
因此同一个微分算子 $\Delta$ 在不同自伴边界域上选出的 调和 代表元 不同：绝对 边界条件给普通 de Rham 上同调类的唯一 调和 代表，相对 边界条件给相对上同调类的唯一 调和 代表。再由 Hodge 星得到
\begin{equation}\label{eq:math-dg-poincare-lefschetz-hodge}
 *:\mathcal H_A^k(M)\xrightarrow{\ \simeq\ }
 \mathcal H_R^{n-k}(M),
\end{equation}
这是 Poincar\'e--Lefschetz 对偶在 调和形式 上的解析实现。

\begin{derivation}[边界 Green 公式如何固定两类自伴域]
对 $\alpha\wedge *\beta$ 作外微分再用 Stokes 定理，得到含边界项的 Green 公式。边界项耦合 $\alpha$ 切向与 $\beta$ 法向数据，因此要求 $\alpha,\beta$ 同时满足 绝对 或 相对 边界条件才能消去边界项，使 $\dd$ 与 $\delta$ 互为 Hilbert 伴随。对 拉普拉斯算子 再做一次分部积分，第二组边界条件消去剩余项，得到非负能量公式。

先把 $\alpha\wedge *\beta$ 作外微分。由 Leibniz 法则，
\begin{equation*}
 \dd(\alpha\wedge *\beta)
 =\dd\alpha\wedge *\beta
 +(-1)^{k-1}\alpha\wedge\dd(*\beta).
\end{equation*}
使用余微分定义把第二项改写成 $(\alpha,\delta\beta)$ 的体积分，再用 Stokes 定理，就得到\eqnrefs{eq:math-dg-hodge-green-boundary}。边界项只耦合 $\alpha$ 的切向数据与 $\beta$ 的法向数据，因此若两边都取 绝对 域，$\iota_{\nu}\beta=0$；若两边都取 相对 域，$j^*\alpha=0$。这正是 $\dd$ 与 $\delta$ 在相应闭域上互为 Hilbert 伴随的原因。

对拉普拉斯算子还要再对 $\dd\omega$ 或 $\delta\omega$ 做一次分部积分。绝对域中的第二条件
$
\iota_{\nu}\dd\omega=0
$
与 相对 域中的第二条件
$
j^*\delta\omega=0
$
分别消去第二次分部积分产生的剩余边界项，所以得到\eqnrefs{eq:math-dg-boundary-hodge-energy}。

若 $\Delta_{A/R,k}\omega=0$，右端是两个非负范数之和，只能同时为零：
\begin{equation*}
 \dd\omega=0,
 \qquad
 \delta\omega=0.
\end{equation*}
因此边界条件并没有修改"调和 $\Rightarrow$ 闭+coclosed"这一局部解析事实；它修改的是允许进入算子定义域的边界数据，从而修改零模空间所实现的上同调理论。

最后检查 Hodge 星。边界上的切--法分解满足
\begin{equation*}
 \iota_{\nu}(*\omega)
 \quad\Longleftrightarrow\quad
 *(\nu^\flat\wedge\omega),
\end{equation*}
差别只含由次数和维数决定的固定符号；同理 $*\dd$ 与 $\delta*$ 也只差固定符号。因此 绝对 的两条条件经 $*$ 恰变成 相对 的两条条件，得到\eqnrefs{eq:math-dg-poincare-lefschetz-hodge}。具体地，若 $\omega$ 满足 $\iota_{\nu}\omega=0$ 与 $\iota_{\nu}\dd\omega=0$（绝对），则 $*\omega$ 满足 $j^*(*\omega)=0$ 与 $j^*\delta(*\omega)=0$（相对），符号由 $(-1)^{k(n-k)}$ 与定向约定决定。

典型例子。
\begin{itemize}
\item 紧致无边界流形：$\partial M=\varnothing$，绝对 与 相对 重合，Green 公式无边界项，Hodge 分解给 $\mathcal H^k(M)\simeq H^k_{\mathrm{dR}}(M)$。
\item 紧致带边流形 $M$：绝对 给 $H^k_{\mathrm{dR}}(M)$，相对 给 $H^k_{\mathrm{dR}}(M,\partial M)$，二者通过 Hodge 星实现 Poincar\'e--Lefschetz 对偶。
\item 区间 $[0,1]$：绝对 边界条件对应 Neumann（法向导数为零），相对 对应 Dirichlet（函数值为零）。调和 $0$-形式：Neumann 给常数（$H^0([0,1])\cong\RR$），Dirichlet 给零（$H^0([0,1],\partial[0,1])=0$）。
\item 圆柱 $M=S^1\times[0,1]$：绝对 在两端固定法向（沿 $[0,1]$ 方向），相对 固定切向（沿 $S^1$ 方向）。调和 $1$-形式空间维数不同，反映绝对与相对上同调的差异。
\item 半球面 $S^n_+$：绝对 与 相对 通过 Hodge 星联系，给出 $H^k(S^n_+)\cong H^k_{\mathrm{dR}}(S^n)$ 与 $H^k(S^n_+,\partial S^n_+)\cong H^{n-k}_{\mathrm{dR}}(S^n)$ 的对偶。
\item Möbius 带：非定向流形上 Hodge 理论需可定向二重覆叠，两类边界条件在二重覆叠上对称分布，反映非定向性的全局拓扑。
\item 单位圆盘 $D^2$：绝对 给 $H^1(D^2)=0$（闭形式恰当），相对 给 $H^1(D^2,\partial D^2)\cong\RR$（相对上同调非平凡），Hodge 星实现 $H^0(D^2)\cong H^2(D^2,\partial D^2)$ 的对偶。
\end{itemize}

两类边界条件对应物理中不同的场边界行为：绝对（Neumann 型）对应"法向通量为零"（如绝缘边界），相对（Dirichlet 型）对应"切向场为零"（如固定边界）。在规范理论中，绝对 边界条件保持规范对称，相对 边界条件固定规范。在热方程方法中，两类边界条件对应不同的热核渐近展开，其系数含边界曲率项（如 Gilkey--Greiner 公式），给出带边流形的指标定理与局部指标密度。Atiyah--Patodi--Singer $\eta$ 算子进一步引入非局部边界条件，把 eta 不变量纳入指标公式，是带边流形指标定理的现代形式。
\end{derivation}

\section{Killing 流与 Hodge 算子的相容性}

Killing 流保持 $g$。因为它从恒等映射连续生成，在局部还保持既定定向，所以同时保持 $\operatorname{vol}_g$。由 Hodge 星的刻画\eqnrefs{eq:math-dg-hodge-star-characterization}可知
\[
\Phi_t^*(*\alpha)=*(\Phi_t^*\alpha).
\]
对 $t=0$ 求导得到
\begin{equation*}[\mathcal L_X,*]=0.
\end{equation*}
再结合 $[\mathcal L_X,\dd]=0$，立即得到
\begin{equation*}[\mathcal L_X,\delta]=0,
\qquad
[\mathcal L_X,\Delta]=0.
\end{equation*}
因此 Hodge 拉普拉斯算子 的每个特征空间都在 Killing 对称下保持。几何对称与谱分解并不是两套独立技巧：它们可以在同一个不变子空间分解中同时使用。